% =============================================================================
% Chapter 13: Model Serving — ML Systems Lecture Slides
% =============================================================================
% IMAGES NEEDED (SVGs converted to PDF):
%   - serving-inversion.pdf          - inference-pipeline.pdf
%   - batching-comparison.pdf        - queuing-latency.pdf
%   - ttft-tpot-timeline.pdf         - continuous-vs-static-batching.pdf
%   - training-serving-skew.pdf
% =============================================================================
\documentclass[aspectratio=169, 12pt]{beamer}
\usepackage{../../assets/beamerthememlsys}

\mlsyssetup{
  volume       = {Volume I},
  chapter      = {Chapter 13},
  logo         = {../../assets/img/logo-mlsysbook.png},
  instlogo     = {../../assets/img/logo-harvard.png},
  chaptertitle = {Model Serving},
}

% --- Fonts ---
\usepackage{fontspec}
\setsansfont{Helvetica Neue}[
  BoldFont={Helvetica Neue Bold},
  ItalicFont={Helvetica Neue Italic},
  BoldItalicFont={Helvetica Neue Bold Italic},
]
\setmonofont{JetBrains Mono}[Scale=0.85]

% --- Packages ---
\usepackage{booktabs}
\usepackage{amsmath}

% --- Image paths ---
\graphicspath{
  {images/}
}

% --- Chapter-specific macros ---
\newcommand{\DAM}{D\raisebox{0.08em}{\tiny$\bullet$}A\raisebox{0.08em}{\tiny$\bullet$}M}

% --- Helper: safe image include ---
\newcommand{\safeimg}[2][width=\textwidth,keepaspectratio]{%
  \IfFileExists{images/#2}{\includegraphics[#1]{#2}}{%
    \IfFileExists{#2}{\includegraphics[#1]{#2}}{%
      \fbox{\parbox[c][2.5cm][c]{0.85\linewidth}{\centering\footnotesize\textcolor{midgray}{[Missing image]}}}%
    }%
  }%
}

% --- Section count for navigation ---
\setcounter{mlsystotalsections}{9}

\title{Model Serving}
\author{Vijay Janapa Reddi}
\institute{Harvard University}
\date{}

\begin{document}

% =============================================================================
% TITLE SLIDE
% =============================================================================
\mlsystitle{Model Serving}{Inverting Every Training Priority}{cover_model_serving.png}

% =============================================================================
% LEARNING OBJECTIVES
% =============================================================================
\begin{frame}{Learning Objectives}
\note{
% -- LINK: What prior concept connects to this slide
Students spent 12 chapters optimizing throughput. This slide frames the
inversion: every training priority is about to flip.

% -- NARRATE: What to SAY while showing this slide
Read each objective aloud, pausing on ``latency budget'' and ``queuing
theory'' --- these are the new quantitative anchors replacing samples/hour.
ANALOGY: ``Training is a factory running 24/7. Serving is an ER --- every
patient has a deadline.''

% -- ENGAGE: Specific question, prediction, or task for THIS slide
Ask: ``How many of you have deployed a model to production? What was your
biggest surprise?'' Cold-call one responder.

% -- WARN: What students will get wrong on THIS topic
Common error: students assume serving is just calling model.predict().
Correct framing: serving is a six-stage pipeline where the model is one
stage consuming less than 50\% of the budget.

% -- FLEX: [CORE] This slide is essential --- do not skip.
IF AHEAD: Ask students to predict which objective will be hardest.
IF SHORT: Read objectives without discussion; move to content.
}

\footnotesize
\begin{enumerate}\setlength\itemsep{0pt}
  \item Explain the \textbf{serving inversion} from throughput to latency minimization
  \item Decompose request latency using the \textbf{latency budget framework}
  \item Apply \textbf{queuing theory} (Little's Law, M/M/1) to capacity planning
  \item Identify sources of \textbf{training-serving skew} and prevention strategies
  \item Select batching strategies based on \textbf{traffic patterns} and latency constraints
  \item Evaluate \textbf{KV-cache}, \textbf{TTFT/TPOT}, and \textbf{continuous batching} for LLM serving
  \item Evaluate deployment tradeoffs across precision, runtime, and cost
\end{enumerate}

\end{frame}

\begin{frame}{Visual Language}
\note{
% -- NARRATE: What to SAY while showing this slide
Point to each card: ``Blue is compute --- GPU forward passes. Green is data
flow and memory. Orange is routing and scheduling. Red flags bottlenecks
and cost.'' In serving diagrams, you will see orange load balancers feeding
blue inference runners, with red marking decode bottlenecks.

% -- FLEX: [OPTIONAL] Skip if students already know the palette from earlier chapters.
IF SHORT: Say ``same colors as always'' and advance.
}

\small
Throughout this course, colors carry meaning:

\vspace{0.3cm}
\begin{columns}[T]
  \begin{column}{0.45\textwidth}
    \begin{mlsyscard}{computestroke}
    \textbf{Blue} --- Compute / Processing\\
    {\footnotesize GPU ops, forward/backward pass, inference}
    \end{mlsyscard}
    \vspace{0.15cm}
    \begin{mlsyscard}{datastroke}
    \textbf{Green} --- Data / Memory\\
    {\footnotesize Data flow, caches, healthy paths}
    \end{mlsyscard}
  \end{column}
  \begin{column}{0.45\textwidth}
    \begin{mlsyscard}{routingstroke}
    \textbf{Orange} --- Routing / Scheduling\\
    {\footnotesize Load balancers, batch windows}
    \end{mlsyscard}
    \vspace{0.15cm}
    \begin{mlsyscard}{errorstroke}
    \textbf{Red} --- Error / Cost / Bottleneck\\
    {\footnotesize Loss, decode phase, waste}
    \end{mlsyscard}
  \end{column}
\end{columns}
\end{frame}


% =============================================================================
\section{The Serving Inversion}
% =============================================================================

\begin{frame}{Why Serving Is Different}
\note{
% -- LINK: What prior concept connects to this slide
In every prior chapter, success meant saturating the GPU. This slide
reveals why that strategy fails in production serving.

% -- NARRATE: What to SAY while showing this slide
Point to the diagram: ``Left side is training --- maximize throughput,
saturate hardware. Right side is serving --- minimize latency, maintain
headroom.'' Walk through the DAM inversion: Data shifts from Volume to
Freshness, Algorithm from Mutable to Frozen, Machine from Saturation to
Headroom.
ANALOGY: ``Training is a freight train --- pack it full. Serving is an
ambulance --- it must always be ready to go.''

% -- ENGAGE: Specific question, prediction, or task for THIS slide
Ask: ``If training saturates GPUs at 95\%, why would serving aim for
50\%?'' Give 15 seconds, then cold-call. [Expected: queuing theory ---
high utilization causes latency spikes.]

% -- WARN: What students will get wrong on THIS topic
Common error: students think serving is just training with batch size 1.
Correct framing: serving inverts the optimization objective itself ---
latency replaces throughput as the primary metric.

% -- FLEX: [CORE] This slide is essential --- do not skip.
IF AHEAD: ``What happens to cost when you run GPUs at 50\% instead of 95\%?''
IF SHORT: Show diagram, state the inversion, skip the DAM walkthrough.
}

% --- Full-width diagram ---
\centering
\safeimg[width=0.88\textwidth,height=4.8cm,keepaspectratio]{serving-inversion.pdf}

\vspace{0.1cm}
\mlsysinsight{Key:}{Training optimizes samples/hour. Serving optimizes ms/request. Same GPU, opposite goals.}

\end{frame}

\begin{frame}{The D\raisebox{0.04em}{\tiny$\bullet$}A\raisebox{0.04em}{\tiny$\bullet$}M Inversion}
\note{
% -- LINK: What prior concept connects to this slide
The previous slide showed the inversion visually. This slide formalizes
it along the three DAM axes students learned in Ch1.

% -- NARRATE: What to SAY while showing this slide
Walk down the table row by row: ``Data: training ingests billions of
samples; serving handles one request at a time --- freshness replaces
volume. Algorithm: training runs backprop; serving is forward-only ---
no optimizer state needed. Machine: training saturates at 95\%; serving
holds headroom at 40--60\% to absorb traffic spikes.'' End on the Iron
Law row: ``The dominant term flips from compute to latency overhead.''

% -- ENGAGE: Specific question, prediction, or task for THIS slide
Before showing the bottom card, ask: ``If both training and serving use
the same GPU, why does the Iron Law's dominant term change?''
[Expected: serving processes one request at low arithmetic intensity,
making the latency/overhead term dominate.]

% -- WARN: What students will get wrong on THIS topic
Common error: students think removing backprop makes serving trivially
easy. Correct framing: removing backprop frees memory but exposes the
latency term that training's large batches amortized away.

% -- FLEX: [CORE] This slide is essential --- do not skip.
IF AHEAD: ``What happens to the Machine row during a traffic spike?''
IF SHORT: Cover Data and Machine rows; skip Algorithm row.
}

\scriptsize
\renewcommand{\arraystretch}{1.1}
\begin{tabular}{@{}llll@{}}
  \toprule
  \textbf{\DAM{} Axis} & \textbf{Training} & \textbf{Serving} & \textbf{Implication} \\
  \midrule
  \textcolor{datastroke}{\textbf{Data}} & Volume (billions) & Freshness (one now) & Cache vs.\ stream \\
  \textcolor{computestroke}{\textbf{Algorithm}} & Mutable (backprop) & Frozen (forward only) & No optimizer state \\
  \textcolor{routingstroke}{\textbf{Machine}} & Saturate (95--100\%) & Headroom (40--60\%) & Absorb traffic spikes \\
  \midrule
  \textbf{Iron Law} & Compute term dominates & Latency term dominates & $L_{\text{lat}}$ is the constraint \\
  \textbf{Metric} & Samples/hour & ms/request (p99) & Tail latency matters \\
  \bottomrule
\end{tabular}

\vspace{0.1cm}
\begin{mlsyscard}{crimson}
{\footnotesize \textbf{The Serving Inversion:} Models that train beautifully often serve poorly. Batch-heavy training optimizations are ill-suited for bursty, latency-critical production traffic.}
\end{mlsyscard}

\end{frame}

\begin{frame}{Static vs.\ Dynamic Inference}
\note{
% -- LINK: What prior concept connects to this slide
The DAM inversion showed that serving prioritizes freshness over volume.
This slide presents the first design decision that follows: when to
compute predictions.

% -- NARRATE: What to SAY while showing this slide
Point to the green card: ``Static inference pre-computes overnight ---
10,000 photos times 5 ms equals 50 seconds total. Zero runtime latency,
but it cannot handle novel inputs.'' Then the red card: ``Dynamic
inference computes on demand under a 100 ms budget. Flexible but
expensive.'' Finish with the hybrid insight at the bottom.

% -- ENGAGE: Specific question, prediction, or task for THIS slide
Ask: ``A search engine autocomplete --- static or dynamic?''
[Expected: hybrid --- common queries are cached, novel queries computed
on demand.]

% -- WARN: What students will get wrong on THIS topic
Common error: students dismiss static inference as outdated. Correct
framing: recommendation systems pre-compute candidate sets for millions
of users nightly; only the final ranking is dynamic.

% -- FLEX: [OPTIONAL] This slide provides context but is not load-bearing.
IF SHORT: Cover the two cards quickly and move on to Server Anatomy.
IF AHEAD: ``What determines the boundary between cached and dynamic?''
}

\footnotesize
\begin{columns}[T]
  \begin{column}{0.48\textwidth}
    \begin{mlsyscard}{datastroke}
    \textbf{Static (Batch) Inference}\\[0.05cm]
    {\scriptsize
    Pre-compute predictions for anticipated inputs.\\
    10,000 photos $\times$ 5 ms $\approx$ 50 s total.\\
    \textbf{Pro:} Eliminates inference latency.\\
    \textbf{Con:} Cannot handle novel inputs.}
    \end{mlsyscard}
  \end{column}
  \begin{column}{0.48\textwidth}
    \begin{mlsyscard}{errorstroke}
    \textbf{Dynamic (Real-time) Inference}\\[0.05cm]
    {\scriptsize
    Compute predictions on demand.\\
    Each request: full pipeline under 100 ms budget.\\
    \textbf{Pro:} Handles any input, reflects updates.\\
    \textbf{Con:} Strict latency, higher cost.}
    \end{mlsyscard}
  \end{column}
\end{columns}

\vspace{0.1cm}
\mlsysconcept{Production:}{Most systems use \textbf{hybrid} --- cached results for common queries, dynamic inference for novel inputs.}

\end{frame}

% =============================================================================
\section{Server Anatomy}
% =============================================================================

\begin{frame}{The Inference Pipeline}
\note{[3 min] Walk through the full pipeline from client to GPU and back.
Key insight: the neural network is only ONE stage. Preprocessing often
consumes 45--70\% of total latency when inference is optimized.
Ask: ``Which stage do you think is the bottleneck?''}

% --- Full-width diagram ---
\centering
\safeimg[width=0.92\textwidth,height=4.5cm,keepaspectratio]{inference-pipeline.pdf}

\vspace{0.1cm}
\mlsysalert{Surprise:}{When inference is fast (5 ms), \textbf{preprocessing dominates} (45--70\% of total latency). Optimize the pipeline, not just the model.}

\end{frame}

\begin{frame}{Inference Server Components}
\note{[3 min] Six stages of the inference server pipeline. Each stage isolates
a concern. The request queue absorbs bursty traffic; the dynamic batcher
forms efficient GPU batches. This is NOT a simple model.predict() call.
Common error: thinking ``serving = calling forward().''}

\scriptsize
\renewcommand{\arraystretch}{1.05}
\begin{tabular}{@{}lp{4cm}p{4cm}@{}}
  \toprule
  \textbf{Component} & \textbf{Purpose} & \textbf{Key Concern} \\
  \midrule
  \textcolor{routingstroke}{\textbf{Load Balancer}} & Route requests to replicas & Least-connections, latency-aware \\
  \textcolor{routingstroke}{\textbf{Network Ingress}} & HTTP/gRPC, TLS termination & Serialization overhead \\
  \textcolor{datastroke}{\textbf{Request Queue}} & Absorb bursty traffic (FIFO) & Queue depth vs.\ latency \\
  \textcolor{computestroke}{\textbf{Dynamic Batcher}} & Group requests by window & Batch size vs.\ wait time \\
  \textcolor{computestroke}{\textbf{Inference Runner}} & Execute model on accelerator & GPU utilization \\
  \textbf{Postprocessor} & Format response (softmax) & Minimal overhead \\
  \bottomrule
\end{tabular}

\vspace{0.1cm}
\mlsysconcept{Key:}{The inference server is a \textbf{high-performance scheduler}, not a model wrapper. Dynamic batching improves GPU utilization by up to 70\%.}

\end{frame}

% --- ACTIVE LEARNING 1: Predict ---
\begin{frame}{Predict: Where Does Time Go?}
\note{[2 min] Prediction exercise before revealing the latency budget.
Give students 60 seconds. Most will guess ``inference'' dominates.
The reveal will surprise them.}

\centering
\vspace{0.6cm}
{\Large\bfseries Think--Write--Share}

\vspace{0.4cm}
{\large A ResNet-50 model runs inference in 5 ms on a GPU.\\
The end-to-end serving latency is 10.1 ms.\\[0.2cm]
\alert{Where does the other 5.1 ms go?}}

\vspace{0.3cm}
{\normalsize Rank the stages: preprocessing, serialization, transfer, postprocessing.}

\vspace{0.3cm}
{\small\textcolor{lightgray}{(60 seconds --- then compare with a neighbor)}}

\end{frame}

% =============================================================================
\section{Latency Budget}
% =============================================================================

\begin{frame}{The Latency Budget Framework}
\note{[3 min] Reveal after prediction exercise. Walk through the budget.
Key surprise: preprocessing (resize, normalize) takes 4.5 ms --- almost as
much as inference itself. The model gets less than 50\% of the total budget.
Common error: optimizing only the model and ignoring preprocessing.}

\footnotesize
\textbf{ResNet-50 inference: Where every millisecond goes}

\vspace{0.1cm}
\renewcommand{\arraystretch}{1.05}
{\scriptsize
\begin{tabular}{@{}llrr@{}}
  \toprule
  \textbf{Stage} & \textbf{Operation} & \textbf{Time} & \textbf{\% Total} \\
  \midrule
  \textcolor{routingstroke}{Deserialize} & Parse JSON, decode payload & 0.2 ms & 2\% \\
  \textcolor{routingstroke}{Preprocess} & Resize, normalize, to-tensor & \textbf{4.5 ms} & \textbf{45\%} \\
  \textcolor{computestroke}{CPU$\to$GPU} & Transfer tensor via PCIe & 0.2 ms & 2\% \\
  \textcolor{computestroke}{Inference} & Neural network forward pass & \textbf{5.0 ms} & \textbf{50\%} \\
  \textcolor{datastroke}{Postprocess} & Softmax, format response & 0.1 ms & 1\% \\
  \midrule
  & \textbf{Total} & \textbf{10.1 ms} & \textbf{100\%} \\
  \bottomrule
\end{tabular}
}

\vspace{0.1cm}
\mlsysalert{Amdahl's Law:}{The model gets $<$50\% of the budget. A 2$\times$ model speedup reduces total latency by only 25\%. Optimize the \emph{pipeline}, not just the model.}

\end{frame}

\begin{frame}{The Cost of Latency}
\note{[2 min] Quantify the dollar cost of latency constraints.
Key insight: halving latency (from 10ms to 5ms) can 4x the cost.
This is why latency requirements must be business-justified.}

\footnotesize
\textbf{GPU rental: \$4/hour. How much does each millisecond cost?}

\vspace{0.1cm}
\renewcommand{\arraystretch}{1.1}
{\scriptsize
\begin{tabular}{@{}lrrr@{}}
  \toprule
  \textbf{Scenario} & \textbf{Latency} & \textbf{Throughput} & \textbf{Cost/M queries} \\
  \midrule
  \textbf{A} (batch=1, low latency) & 5 ms & 200 req/s & \textbf{\$5.56} \\
  \textbf{B} (batch=8, high throughput) & 10 ms & 800 req/s & \textbf{\$1.39} \\
  \bottomrule
\end{tabular}
}

\vspace{0.1cm}
\mlsysalert{Trade-off:}{Reducing latency from 10 ms to 5 ms increases the hardware bill by \textbf{$\approx$300\%}. Latency requirements must be business-justified.}

\end{frame}

% =============================================================================
\section{Queuing Theory}
% =============================================================================

\mlsysfocus{Little's Law}{%
$L = \lambda \cdot W$\\[0.3cm]
{\normalsize Average queue length = Arrival rate $\times$ Average wait time}%
}

\begin{frame}{The Utilization-Latency Curve}
\note{[3 min] Walk through the M/M/1 queuing curve. Key insight: latency is
NOT linear with utilization. At 70\%, latency is 3.3x service time. At 90\%,
it is 10x. This is why production systems run at 40-60\% utilization.
Ask: ``Your manager wants to run GPUs at 90\% to save money. What do you say?''}

% --- Full-width diagram ---
\centering
\safeimg[width=0.88\textwidth,height=4.5cm,keepaspectratio]{queuing-latency.pdf}

\vspace{0.1cm}
\mlsysalert{The Knee:}{At $\sim$70\% utilization, latency starts exploding. Moving 70\%$\to$90\% saves $\sim$22\% cost but \textbf{triples} latency.}

\end{frame}

\begin{frame}{Quick Check}
\note{[0.5 min] Micro-retrieval. Give students 30 seconds to recall.
This distributed practice improves retention (Roediger \& Karpicke, 2006).}

\centering
\Large\bfseries Quick check --- no peeking.\\[0.5cm]
\normalsize At what GPU utilization does queuing latency start to explode?\\[0.3cm]
{\small 30 seconds --- then we continue.}
\end{frame}



\begin{frame}{M/M/1 Wait Time}
\note{[3 min] The governing equation. Walk through each variable.
At 80\% utilization with 5ms service time: wait = 20ms, total = 25ms.
The p99 is approximately 4.6x the mean: 115ms.
This single equation explains why production systems cannot run hot.}

\footnotesize
$$W_q = \frac{\rho}{1 - \rho} \cdot \mu^{-1} \qquad \text{where } \rho = \frac{\lambda}{\mu}$$

\vspace{0.05cm}
{\scriptsize
\begin{tabular}{@{}ll@{}}
  $W_q$ --- Average wait time in queue &
  $\rho$ --- Utilization ($\lambda / \mu$) \\
  $\lambda$ --- Arrival rate (req/s) &
  $\mu$ --- Service rate (1 / service time) \\
\end{tabular}
}

\vspace{0.05cm}
\renewcommand{\arraystretch}{1.0}
{\scriptsize
\begin{tabular}{@{}rrrrr@{}}
  \toprule
  \textbf{Utilization} & \textbf{Wait} & \textbf{Total} & \textbf{p99 ($\approx$4.6$\times$)} & \textbf{Assessment} \\
  \midrule
  50\% & 1$\times$ & 2$\times$ & 9.2$\times$ & \textcolor{datastroke}{Safe} \\
  70\% & 2.3$\times$ & 3.3$\times$ & 15.3$\times$ & \textcolor{routingstroke}{Caution} \\
  80\% & 4$\times$ & 5$\times$ & 23$\times$ & \textcolor{errorstroke}{Danger} \\
  90\% & 9$\times$ & 10$\times$ & 46$\times$ & \textcolor{errorstroke}{Critical} \\
  \bottomrule
\end{tabular}
}

\end{frame}

\begin{frame}{War Story: The Tail at Scale}
\note{
% -- LINK: What prior concept connects to this slide
The M/M/1 model showed why high utilization explodes latency for one
server. This slide extends the insight to fan-out systems with hundreds
of servers --- the production reality at Google, Meta, and Amazon.

% -- NARRATE: What to SAY while showing this slide
``Dean and Barroso, 2013: a single server is slow 1\% of the time. With
100 servers, the probability that at least one is slow is $1 - 0.99^{100}
= 63\%$. With 1,000 servers: 99.99\%. Your p99 at the component level
becomes the median at the system level.'' Point to the table: ``At 100
servers, most requests hit at least one straggler. Hedged requests ---
sending the same query to two servers and taking the first response ---
reduce tail latency by 75\%.''
ANALOGY: ``One slow cashier at a supermarket is an annoyance. One slow
cashier across 100 checkout lanes means 63\% of shoppers wait.''

% -- ENGAGE: Specific question, prediction, or task for THIS slide
Ask: ``Your search query fans out to 1,000 shards. Each shard has p99 =
10 ms. What is the system-level p99?''
[Expected: much worse than 10 ms --- the slowest shard dominates.]

% -- WARN: What students will get wrong on THIS topic
Common error: students design for single-server tail latency and forget
fan-out amplification. Correct framing: at scale, component-level
percentiles are nearly meaningless --- system-level tail is what users
experience.

% -- FLEX: [CORE] The Tail at Scale insight is essential for production systems
[CORE] Connects queuing theory to real-world distributed serving.
IF AHEAD: ``How do hedged requests interact with utilization?''
IF SHORT: State the 63\% number and the hedging mitigation.
}

\footnotesize
\textbf{Dean \& Barroso (2013): ``The Tail at Scale''}

\vspace{0.1cm}
\begin{columns}[T]
  \begin{column}{0.52\textwidth}
    {\scriptsize
    One server slow 1\% of the time.\\
    A request touching $N$ servers:
    $$P(\text{$\geq$1 slow}) = 1 - 0.99^N$$
    }

    \vspace{-0.1cm}
    {\scriptsize
    \renewcommand{\arraystretch}{1.0}
    \begin{tabular}{@{}rr@{}}
      \toprule
      \textbf{Servers} & \textbf{P($\geq$1 slow)} \\
      \midrule
      1    & 1\% \\
      10   & 10\% \\
      100  & \textbf{63\%} \\
      1000 & 99.99\% \\
      \bottomrule
    \end{tabular}
    }
  \end{column}
  \begin{column}{0.45\textwidth}
    \begin{mlsyscard}{routingstroke}
    \textbf{Mitigations}\\[0.05cm]
    {\scriptsize
    \textbf{Hedged requests:} Send to 2 replicas, take first response. Tail $\downarrow$ 75\%.\\[0.05cm]
    \textbf{Tied requests:} Second server cancels if first responds. Near-zero overhead.\\[0.05cm]
    \textbf{Micro-partitioning:} Fine-grained shards for load balance.}
    \end{mlsyscard}
  \end{column}
\end{columns}

\vspace{0.05cm}
\mlsysalert{At scale:}{Component p99 becomes the system \textbf{median}. Design for the tail, not the average.}

\end{frame}

% --- ACTIVE LEARNING 2: Exercise ---
\begin{frame}{Your Turn: Capacity Planning}
\note{[3 min] Give students 90 seconds. Answer: At 200 req/s with 5ms
service time, rho = 200 * 0.005 = 1.0 --- system is AT capacity.
Average wait is infinite. Need to either reduce service time or add servers.
PEER INSTRUCTION PROTOCOL:
1. Present problem (30s). 2. Individual vote (60s).
3. If 30-70\% correct: pair discussion (90s), then re-vote.
4. Reveal and explain.}

\footnotesize
\begin{columns}[T]
  \begin{column}{0.55\textwidth}
    {\small\bfseries Capacity Planning Exercise}

    \vspace{0.1cm}
    A serving system has:
    \begin{itemize}\setlength\itemsep{0pt}
      \item Service time: 5 ms per request
      \item Expected traffic: 200 req/s
      \item SLO: p99 latency $<$ 100 ms
    \end{itemize}

    \vspace{0.1cm}
    \textbf{Calculate:} (1) utilization $\rho$, (2) can one server meet the SLO?, (3) how many servers needed?

    \vspace{0.1cm}
    {\scriptsize\textcolor{midgray}{(90 seconds --- then compare with a neighbor)}}
  \end{column}
  \begin{column}{0.42\textwidth}
    \pause
    \begin{mlsyscard}{computestroke}
    \textbf{Solution:}\\[0.05cm]
    {\scriptsize
    $\rho = 200 \times 0.005 = 1.0$\\[0.03cm]
    \textcolor{errorstroke}{\textbf{At 100\%!}} Infinite queue.\\[0.05cm]
    Need $\rho < 0.7$ for p99 target.\\
    Min: $\lceil 200 / (200 \times 0.7) \rceil$ = \textbf{2 servers}.\\[0.05cm]
    With 2: $\rho = 0.5$, p99 $\approx$ 46 ms \textcolor{datastroke}{\checkmark}
    }
    \end{mlsyscard}
  \end{column}
\end{columns}

\end{frame}

% =============================================================================
\section{Batching Strategies}
% =============================================================================

\begin{frame}{Why Batching Helps}
\note{[3 min] Full-width timeline comparing batch-1 vs batch-8.
Key numbers: 200 vs 879 img/s (4.4x), cost drops from \$5.56 to \$1.26
per million queries. GPU goes from 50\% idle to 65\% utilized.
Ask: ``Why not batch=1024?'' (answer: latency SLO violation)}

% --- Full-width diagram ---
\centering
\safeimg[width=0.88\textwidth,height=4.2cm,keepaspectratio]{batching-comparison.pdf}

\vspace{0.05cm}
\mlsysconcept{Trade-off:}{Larger batches improve throughput and cost, but increase latency. The \textbf{SLO} sets the upper bound.}

\end{frame}

\begin{frame}{Static vs.\ Dynamic Batching}
\note{[2 min] Two batching strategies. Static: fixed batch size, wait for N
requests. Dynamic: time-window-based, fires when window expires OR batch is
full. Dynamic is strictly better for variable traffic but adds complexity.}

\small
\begin{columns}[T]
  \begin{column}{0.48\textwidth}
    \begin{mlsyscard}{routingstroke}
    \textbf{Static Batching}\\[0.1cm]
    {\footnotesize
    Fixed batch size (e.g., 8).\\
    Wait until batch is full.\\[0.1cm]
    \textbf{Pro:} Simple implementation.\\
    \textbf{Con:} High latency during low traffic (waiting for requests to fill the batch).}
    \end{mlsyscard}
  \end{column}
  \begin{column}{0.48\textwidth}
    \begin{mlsyscard}{computestroke}
    \textbf{Dynamic Batching}\\[0.1cm]
    {\footnotesize
    Time window (e.g., 5 ms) + max size.\\
    Fire when \emph{either} condition met.\\[0.1cm]
    \textbf{Pro:} Adapts to traffic patterns.\\
    \textbf{Con:} Tuning window size is critical (too small $\to$ small batches; too large $\to$ SLO violation).}
    \end{mlsyscard}
  \end{column}
\end{columns}

\vspace{0.15cm}
\mlsysinsight{Production pattern:}{Dynamic batching with adaptive windows is the default in NVIDIA Triton, TF Serving, and TorchServe.}

\end{frame}

\begin{frame}{Traffic Patterns and Batching Strategy}
\note{[2 min] Match batching strategy to traffic pattern. Four patterns from
MLPerf scenarios. The deployment paradigm from Ch2 determines which pattern
applies. If short: just cover Server and SingleStream.}

\scriptsize
\renewcommand{\arraystretch}{1.05}
\begin{tabular}{@{}lp{2.5cm}p{2.2cm}p{3.5cm}@{}}
  \toprule
  \textbf{MLPerf Scenario} & \textbf{Traffic Pattern} & \textbf{Batching} & \textbf{Example} \\
  \midrule
  \textbf{Server} & Poisson arrivals & Dynamic window & Web API, recommendation \\
  \textbf{SingleStream} & One at a time & None (batch=1) & Mobile app, edge device \\
  \textbf{MultiStream} & Synchronized sensors & Fixed batch & Autonomous vehicle cameras \\
  \textbf{Offline} & All data available & Max batch & Nightly batch processing \\
  \bottomrule
\end{tabular}

\vspace{0.1cm}
\mlsysconcept{Key:}{The deployment paradigm determines the traffic pattern. Cloud $\to$ Server. Mobile $\to$ SingleStream. Edge $\to$ MultiStream.}

\end{frame}

% =============================================================================
\section{LLM Serving}
% =============================================================================

\begin{frame}{LLM Serving: Two Phases, Two Bottlenecks}
\note{[3 min] The defining challenge of LLM serving. Walk through the timeline.
Prefill is compute-bound (all tokens in parallel). Decode is memory-bandwidth
bound (one token at a time, read entire model per token).
These are governed by DIFFERENT hardware limits.
Ask: ``Adding more compute cores would help which phase?''}

% --- Full-width diagram ---
\centering
\safeimg[width=0.88\textwidth,height=4.2cm,keepaspectratio]{ttft-tpot-timeline.pdf}

\vspace{0.05cm}
{\footnotesize
\textcolor{computestroke}{\textbf{TTFT}} (Time to First Token): Responsiveness. \quad
\textcolor{errorstroke}{\textbf{TPOT}} (Time Per Output Token): Fluidity.}

\end{frame}

\begin{frame}{The KV Cache: Memory That Grows}
\note{[3 min] The defining memory constraint of LLM serving. KV cache grows
linearly with sequence length AND batch size. For a 70B model, batch 32 at
8K context requires 330 GB --- exceeding model weights (140 GB).
Key numbers: 1.3 MB per token for a 70B model.}

\small
\begin{columns}[T]
  \begin{column}{0.52\textwidth}
    \textbf{KV Cache grows with every token:}

    \vspace{0.15cm}
    {\footnotesize
    \renewcommand{\arraystretch}{1.15}
    \begin{tabular}{@{}lrr@{}}
      \toprule
      \textbf{Scenario} & \textbf{KV Cache} & \textbf{vs.\ Weights} \\
      \midrule
      70B, batch 1, 640 tok & 832 MB & 0.6$\times$ \\
      70B, batch 8, 8K tok & 83 GB & 0.6$\times$ \\
      70B, batch 32, 8K tok & \textbf{330 GB} & \textbf{2.4$\times$} \\
      \bottomrule
    \end{tabular}
    }

    \vspace{0.15cm}
    \mlsysalert{OOM risk:}{At batch 32 with 8K context, KV cache \textbf{exceeds model weights} and overflows a single A100 (80 GB).}
  \end{column}
  \begin{column}{0.44\textwidth}
    \begin{mlsyscard}{datastroke}
    \textbf{Solutions}\\[0.1cm]
    {\footnotesize
    \textbf{PagedAttention:} Fixed-size memory pages, near-zero waste (OS virtual memory concept).\\[0.1cm]
    \textbf{Prefix Caching:} Reuse KV states for shared system prompts.\\[0.1cm]
    \textbf{KV Offloading:} Spill inactive contexts to CPU RAM or SSD.}
    \end{mlsyscard}
  \end{column}
\end{columns}

\end{frame}

\begin{frame}{Continuous Batching}
\note{[3 min] Walk through static vs continuous batching grids.
Static: batch locked until longest request finishes, wasting 40\% compute.
Continuous: slots reused immediately, 95\% utilization, 2-4x throughput.
This is the key innovation enabling practical LLM serving (vLLM, TGI).}

% --- Full-width diagram ---
\centering
\safeimg[width=0.92\textwidth,height=4.8cm,keepaspectratio]{continuous-vs-static-batching.pdf}

\vspace{0.05cm}
\mlsysinsight{vLLM innovation:}{Continuous batching + PagedAttention together achieve 2--4$\times$ throughput over naive static batching.}

\end{frame}

\begin{frame}{Node-Level Optimization}
\note{
% -- LINK: What prior concept connects to this slide
Continuous batching and PagedAttention optimize the software layer.
This slide addresses the hardware layer: how to extract maximum throughput
from a single GPU node before scaling out.

% -- NARRATE: What to SAY while showing this slide
``Three techniques that compound: Multi-Instance GPU (MIG) partitions one
H100 into up to 7 isolated instances --- run different models or batch
sizes on each. Model parallelism splits a model too large for one GPU
across 2--4 GPUs within a node using NVLink at 900 GB/s. Speculative
decoding uses a small draft model to propose tokens, then the large model
verifies in parallel --- 2--3x decode speedup.''
ANALOGY: ``MIG is like partitioning a hard drive. Model parallelism is
like RAID striping. Speculative decoding is like a junior associate
drafting a brief for a senior partner to review.''

% -- ENGAGE: Specific question, prediction, or task for THIS slide
Ask: ``When would you choose MIG over model parallelism?''
[Expected: MIG for multiple small models; model parallelism for one
large model that does not fit in one GPU.]

% -- WARN: What students will get wrong on THIS topic
Common error: students assume model parallelism always helps. Correct
framing: NVLink bandwidth (900 GB/s) is still 4x slower than HBM3
(3.35 TB/s). Parallelism adds communication overhead --- only use it
when the model does not fit.

% -- FLEX: [OPTIONAL] Deepens the single-node picture before scaling out
[OPTIONAL] Supplements the core LLM serving content.
IF AHEAD: ``What is the communication overhead of tensor parallelism?''
IF SHORT: Name the three techniques with one sentence each, move on.
}

\scriptsize
\renewcommand{\arraystretch}{1.05}
\begin{tabular}{@{}lp{3.5cm}p{3.5cm}@{}}
  \toprule
  \textbf{Technique} & \textbf{How It Works} & \textbf{When to Use} \\
  \midrule
  \textbf{MIG} & Partition 1 GPU into 7 instances & Multiple small models / right-size GPU slices \\
  \textbf{Tensor Parallelism} & Split layers across 2--8 GPUs via NVLink (900 GB/s) & Model $>$ 1 GPU memory \\
  \textbf{Speculative Decode} & Small draft model proposes, large model verifies & Decode-bound LLMs (2--3$\times$ speedup) \\
  \bottomrule
\end{tabular}

\vspace{0.15cm}
\mlsysconcept{Key:}{Exhaust single-node optimizations \textbf{before} scaling to multiple nodes. NVLink (900 GB/s) $\gg$ network (400 Gbps = 50 GB/s).}

\end{frame}

% --- Bridge slide: Batching to LLM Serving ---
\begin{frame}{From Fixed to Variable: Why LLMs Change Everything}
\note{
% -- LINK: What prior concept connects to this slide
Students just learned batching strategies for fixed-size inputs (images).
This slide bridges to LLM serving where outputs are variable-length.

% -- NARRATE: What to SAY while showing this slide
``Everything we said about batching assumed fixed-size inputs and outputs.
A ResNet batch of 8 images: 8 identical forward passes, done. An LLM
batch of 8 prompts: each generates a different number of tokens. Request
A finishes in 50 tokens; Request H needs 500. Static batching wastes
90\% of compute waiting for the longest request.''

% -- ENGAGE: Specific question, prediction, or task for THIS slide
Ask: ``If one request in a batch of 8 generates 10x more tokens than
the others, what fraction of GPU time is wasted?''
[Expected: roughly 7/8 of compute is idle for most of the batch duration.]

% -- WARN: What students will get wrong on THIS topic
Common error: students apply image-batching intuition to LLMs. Correct
framing: variable output lengths fundamentally change the batching
calculus --- continuous batching exists to solve this specific problem.

% -- FLEX: [CORE] Essential bridge between batching and LLM serving
[CORE] Without this bridge, students will not understand why continuous
batching is necessary. IF SHORT: State the key contrast in 30 seconds.
}

\centering
\vspace{0.3cm}
\footnotesize
\begin{columns}[T]
  \begin{column}{0.46\textwidth}
    \begin{mlsyscard}{datastroke}
    \textbf{Fixed-Size Inference}\\[0.1cm]
    {\scriptsize
    ResNet: batch of 8 images.\\
    All 8 finish simultaneously.\\
    GPU utilization: \textbf{high}.\\
    Batching strategy: straightforward.}
    \end{mlsyscard}
  \end{column}
  \begin{column}{0.46\textwidth}
    \begin{mlsyscard}{errorstroke}
    \textbf{Variable-Length Generation}\\[0.1cm]
    {\scriptsize
    LLM: batch of 8 prompts.\\
    Request A: 50 tokens. Request H: 500.\\
    GPU wastes \textbf{$\sim$90\%} waiting.\\
    Static batching: catastrophic waste.}
    \end{mlsyscard}
  \end{column}
\end{columns}

\vspace{0.2cm}
\mlsysalert{The problem:}{Variable output lengths break fixed batching. \textbf{Continuous batching} fills slots as requests finish --- the key innovation for practical LLM serving.}

\end{frame}

% --- ACTIVE LEARNING 3: Discussion ---
\begin{frame}{Discussion: The Decode Bottleneck}
\note{[3 min] Turn-and-talk. Students discuss in pairs for 90 seconds.
Key insight: adding more FLOPS does NOT help decode because it is
memory-bandwidth bound. Only faster memory or smaller models help.
Cold-call 2--3 pairs.}

\centering
\vspace{0.5cm}
{\large\bfseries Turn and Talk \textcolor{midgray}{(90 seconds)}}

\vspace{0.4cm}
{\large An LLM generates tokens at 20 tok/s on an H100 (989 TFLOPS).\\
You upgrade to a GPU with 2$\times$ compute but same memory bandwidth.\\[0.2cm]
\alert{Does token generation speed double?}}

\vspace{0.3cm}
{\normalsize Hint: Is decode compute-bound or memory-bandwidth-bound?}

\end{frame}

% --- ACTIVE LEARNING 4: Quantitative Exercise ---
\begin{frame}{Your Turn: Llama-3-70B Decode Throughput}
\note{
% -- LINK: What prior concept connects to this slide
The discussion established that decode is memory-bandwidth bound. This
exercise forces students to calculate the exact throughput limit.

% -- NARRATE: What to SAY while showing this slide
``Llama-3-70B at FP16: 140 GB of weights. The H100 has 3.35 TB/s HBM3
bandwidth. Decode reads the FULL model per token. At batch=1: 140 GB /
3.35 TB/s = 41.8 ms per token = 23.9 tok/s. At batch=8: still 41.8 ms
per token generation step, but you produce 8 tokens per step = 191
tok/s aggregate. The model weights are read once per step regardless of
batch size --- batching amortizes the bandwidth cost.''

% -- ENGAGE: Peer Instruction Protocol
Present problem (30s). Individual calculation (60s). If 30--70\%
correct: pair discussion (90s), re-vote.
[Expected errors: students multiply bandwidth by batch size, or forget
that batch=8 still reads weights only once per step.]

% -- WARN: What students will get wrong on THIS topic
Common error: students think batch=8 requires 8x the bandwidth. Correct
framing: all 8 tokens share the same weight read. Batching increases
arithmetic intensity without increasing bandwidth demand.

% -- FLEX: [CORE] The quantitative centerpiece of LLM serving
[CORE] The 23.9 tok/s number anchors Key Takeaways and cost analysis.
IF AHEAD: ``At what batch size does this become compute-bound?''
IF SHORT: Show the solution directly, skip peer discussion.
}

\footnotesize
\begin{columns}[T]
  \begin{column}{0.55\textwidth}
    {\small\bfseries LLM Decode Throughput}

    \vspace{0.1cm}
    Llama-3-70B at FP16:
    \begin{itemize}\setlength\itemsep{0pt}
      \item Model weights: 70B $\times$ 2 B = \textbf{140 GB}
      \item H100 HBM3 bandwidth: \textbf{3.35 TB/s}
      \item Decode reads full model per token step
    \end{itemize}

    \vspace{0.1cm}
    \textbf{Calculate:}
    \begin{enumerate}\setlength\itemsep{0pt}
      \item Max tok/s at batch=1?
      \item Max tok/s at batch=8?
      \item Why does batching help?
    \end{enumerate}

    {\scriptsize\textcolor{midgray}{(90 seconds --- compare with a neighbor)}}
  \end{column}
  \begin{column}{0.42\textwidth}
    \pause
    \begin{mlsyscard}{computestroke}
    \textbf{Solution:}\\[0.05cm]
    {\scriptsize
    \textbf{Batch=1:}\\
    $140\text{ GB} / 3.35\text{ TB/s} = 41.8$ ms/step\\
    $= $ \textbf{23.9 tok/s}\\[0.08cm]
    \textbf{Batch=8:}\\
    Same 41.8 ms/step, but 8 tokens/step\\
    $= $ \textbf{191 tok/s} aggregate\\[0.08cm]
    \textcolor{datastroke}{\textbf{Why:}} Weights read once per step.\\
    Batch amortizes bandwidth cost.
    }
    \end{mlsyscard}
  \end{column}
\end{columns}

\end{frame}

% =============================================================================
\section{Lifecycle \& Skew}
% =============================================================================

\begin{frame}{Training-Serving Skew}
\note{[3 min] Full-width diagram. Walk through specific examples of skew.
PIL bilinear vs OpenCV bicubic resize alone can cause 0.5-1\% accuracy drop.
The model weights are IDENTICAL. Standard monitoring catches nothing.
This is the most insidious failure mode in serving.}

% --- Full-width diagram ---
\centering
\safeimg[width=0.88\textwidth,height=4.5cm,keepaspectratio]{training-serving-skew.pdf}

\vspace{0.1cm}
\mlsysalert{Silent failure:}{Same weights + different preprocessing = different predictions. 5\% accuracy loss invisible to latency monitoring.}

\end{frame}

\begin{frame}{Cold Start and Model Loading}
\note{[2 min] Cold start is not just loading weights. Graph compilation and
memory allocation often take longer than the data transfer. Cloud cold start
can be seconds to minutes. Mobile: milliseconds. TinyML: microseconds.}

\scriptsize
\begin{columns}[T]
  \begin{column}{0.52\textwidth}
    {\footnotesize\textbf{Cold start = more than weight loading:}}

    \vspace{0.1cm}
    \renewcommand{\arraystretch}{1.0}
    \begin{tabular}{@{}lr@{}}
      \toprule
      \textbf{Phase} & \textbf{Time} \\
      \midrule
      Container startup & 5--30 s \\
      Weight loading (CPU$\to$GPU) & 2--10 s \\
      Graph compilation / JIT & 5--60 s \\
      Memory allocation (KV cache) & 1--5 s \\
      Warmup inference & 1--5 s \\
      \midrule
      \textbf{Total cold start} & \textbf{15--110 s} \\
      \bottomrule
    \end{tabular}
  \end{column}
  \begin{column}{0.45\textwidth}
    \vspace{0.1cm}
    \safeimg[width=\textwidth,height=2.5cm,keepaspectratio]{server-rack-datacenter.jpg}

    \vspace{0.05cm}
    {\scriptsize\textcolor{midgray}{Production serving infrastructure. CC BY 2.0.}}

    \vspace{0.1cm}
    \begin{mlsyscard}{routingstroke}
    {\scriptsize\textbf{Mitigations:} model caching, pre-warming, serialized engines, keep-alive (never scale to zero).}
    \end{mlsyscard}
  \end{column}
\end{columns}

\end{frame}

% =============================================================================
\section{Runtime \& Cost}
% =============================================================================

\begin{frame}{Inference Runtime Selection}
\note{[2 min] The runtime can vary by 10x in performance for the same model.
Walk through the three tiers: framework-native (simplest), ONNX (portable),
TensorRT (fastest but vendor-locked). Each has clear tradeoffs.}

\scriptsize
\renewcommand{\arraystretch}{1.0}
\begin{tabular}{@{}lp{2.2cm}p{2.2cm}p{2.8cm}@{}}
  \toprule
  \textbf{Runtime} & \textbf{Speedup} & \textbf{Pro} & \textbf{Con} \\
  \midrule
  \textbf{Framework-native} (PyTorch, TF) & 1$\times$ (baseline) & Any model works & Training overhead \\
  \textbf{ONNX Runtime} & 2--3$\times$ & Cross-platform & Export step required \\
  \textbf{TensorRT} & 3--5$\times$ & Max GPU perf & NVIDIA-only \\
  \textbf{vLLM / TGI} & 2--4$\times$ (LLM) & Continuous batching & LLM-specific \\
  \bottomrule
\end{tabular}

\vspace{0.1cm}
\mlsysconcept{Key:}{Runtime selection is as impactful as model architecture. TensorRT delivers 3--5$\times$ speedup = 3--5$\times$ fewer GPUs.}

\end{frame}

\begin{frame}{Precision Selection for Serving}
\note{[2 min] Extend the quantization concepts from Ch10 into serving.
INT8 gives ~3x throughput over FP32 with careful calibration.
Key pitfall: calibrate with PRODUCTION traffic, not training data.
FP16/BF16 is the safe default; INT8 requires validation.}

\scriptsize
\renewcommand{\arraystretch}{1.0}
\begin{tabular}{@{}lrrrr@{}}
  \toprule
  \textbf{Precision} & \textbf{Bytes} & \textbf{Throughput} & \textbf{Accuracy Risk} & \textbf{Use Case} \\
  \midrule
  FP32 & 4 B & 1$\times$ & None & Debugging, reference \\
  FP16 / BF16 & 2 B & $\sim$2$\times$ & Minimal & Safe default for serving \\
  INT8 & 1 B & $\sim$3$\times$ & Low--moderate & Cost-sensitive, validated \\
  INT4 & 0.5 B & $\sim$5$\times$ & Moderate--high & Research, edge \\
  \bottomrule
\end{tabular}

\vspace{0.05cm}
\footnotesize
\mlsysalert{Calibration pitfall:}{Calibrate with \textbf{production traffic}, not training data. Distribution mismatch invalidates quantization.}

\end{frame}

\begin{frame}{Serving Economics: Llama-3-8B Case Study}
\note{[3 min] Concrete cost analysis. Walk through the numbers for serving
Llama-3-8B on an H100. Key takeaways: INT8 halves cost, batching halves it
again. Runtime + precision choices directly determine financial viability.
The intelligence deflation trend makes margins tighter every year.}

\footnotesize
\textbf{Serving Llama-3-8B on H100 (\$4/hour)}

\vspace{0.05cm}
\renewcommand{\arraystretch}{1.0}
{\scriptsize
\begin{tabular}{@{}llrrr@{}}
  \toprule
  \textbf{Config} & \textbf{Precision} & \textbf{Throughput} & \textbf{Cost/M tokens} & \textbf{Savings} \\
  \midrule
  Naive (batch=1, PyTorch) & FP16 & 50 tok/s & \$22.22 & --- \\
  + TensorRT & FP16 & 150 tok/s & \$7.41 & 3.0$\times$ \\
  + INT8 quantization & INT8 & 400 tok/s & \$2.78 & 8.0$\times$ \\
  + Dynamic batching (b=32) & INT8 & 1,200 tok/s & \$0.93 & \textbf{24$\times$} \\
  \bottomrule
\end{tabular}
}

\vspace{0.1cm}
\mlsysconcept{Key:}{Runtime (3$\times$) + quantization (2.7$\times$) + batching (3$\times$) = \textbf{24$\times$ cost reduction} on the same GPU.}

\end{frame}

% =============================================================================
\section{Wrap-Up}
% =============================================================================

\begin{frame}{Fallacies}
\note{[2 min] Four common misconceptions with quantitative evidence.}

\small
\textbf{Fallacy:} \textit{Reducing model inference latency proportionally reduces user-perceived latency.}\\
{\footnotesize At 80\% utilization, queuing wait = 20 ms on top of 5 ms inference. A 2.5$\times$ model speedup yields 8.5$\times$ system speedup due to nonlinear queuing effects.}

\vspace{0.1cm}
\textbf{Fallacy:} \textit{Training accuracy guarantees serving accuracy.}\\
{\footnotesize PIL vs.\ OpenCV resize alone shifts accuracy by 0.5--1\%. A model at 95\% validation drops to 90\% in production from preprocessing mismatches --- invisible to latency monitoring.}

\vspace{0.1cm}
\textbf{Fallacy:} \textit{Larger serving batches always improve throughput.}\\
{\footnotesize Beyond the optimal batch size, queuing delay exceeds throughput gains. The SLO sets a hard ceiling: batch=64 may maximize throughput but violate the 50 ms p99 target.}

\end{frame}

\begin{frame}{Pitfalls}
\note{[2 min] Three operational pitfalls.}

\scriptsize
\textbf{Pitfall:} \textit{Running serving infrastructure at high utilization for cost efficiency.}\\
70\%$\to$90\% utilization saves $\sim$22\% cost but \textbf{triples latency}. p99 jumps from $\sim$77 ms to $\sim$230 ms.

\vspace{0.08cm}
\textbf{Pitfall:} \textit{Using average latency to evaluate serving performance.}\\
At 70\% utilization, mean = $\sim$17 ms but p99 = $\sim$77 ms (4.6$\times$ gap). Average dashboards mask 1-in-100 user experience.

\vspace{0.08cm}
\textbf{Pitfall:} \textit{Calibrating quantized models with training data rather than production traffic.}\\
INT8 models calibrated on training data can lose 2--5\% accuracy on production inputs with different dynamic ranges.

\end{frame}


\begin{frame}{Muddiest Point}
\note{[1 min] One-minute paper. Collect responses (physical or digital).
Use responses to open the NEXT lecture with targeted clarification.
Angelo \& Cross (1993) --- powerful diagnostic tool.}

\centering
\Large\bfseries One-minute paper\\[0.5cm]
\normalsize Write down the \textbf{muddiest point} from today's lecture.\\[0.2cm]
{\small What concept was most confusing or unclear?}\\[0.5cm]
{\footnotesize\textcolor{midgray}{(Hand in on your way out --- or submit digitally)}}
\end{frame}

% --- RETRIEVAL PRACTICE ---
\begin{frame}{What Were the Key Ideas?}
\note{[2 min] Retrieval practice. Students write 90 seconds, no notes.}

\centering
\vspace{1.5cm}
{\Large\bfseries Close your notes.}

\vspace{0.8cm}
{\large Write down the \textbf{4 most important concepts} from today.}

\vspace{0.8cm}
{\normalsize\textcolor{midgray}{90 seconds --- no peeking.}}

\end{frame}

\begin{frame}{Key Takeaways}
\note{[2 min] Reveal. Walk through each bullet. Emphasize quantitative anchors:
4.4x batching, 24x optimization stack, 3.3x queuing at 70\%, 5\% skew loss.}

\scriptsize
\begin{itemize}\setlength\itemsep{0pt}
  \item \textbf{Serving inverts training}: Throughput $\to$ latency. Saturation $\to$ headroom. Volume $\to$ freshness.
  \item \textbf{Preprocessing dominates}: When inference is fast (5 ms), preprocessing consumes 45--70\% of total latency. Optimize the pipeline.
  \item \textbf{Queuing theory governs capacity}: At 80\% util, wait = 5$\times$ service time. At 90\%, 10$\times$. The knee at $\sim$70\% sets the ceiling.
  \item \textbf{Batching follows traffic}: Poisson arrivals use dynamic batching; streaming sensors use synchronized batches; mobile eliminates batching.
  \item \textbf{Training-serving skew}: Different preprocessing causes 5\% accuracy loss invisible to standard monitoring. Use identical code paths.
  \item \textbf{LLM serving is memory-BW bound}: Decode reads entire model per token. Continuous batching + PagedAttention = 2--4$\times$ throughput.
  \item \textbf{Runtime + precision = cost}: TensorRT + INT8 + batching = 24$\times$ cost reduction on identical hardware.
\end{itemize}

\end{frame}

\begin{frame}{References}
\note{[1 min] Point students to canonical papers.}

\small
\mlsysref{Kwon+23}{Kwon et al. ``Efficient Memory Management for LLM Serving with PagedAttention.'' SOSP 2023.}
\mlsysref{Yu+22}{Yu et al. ``Orca: A Distributed Serving System for Transformer-Based Models.'' OSDI 2022.}
\mlsysref{Barroso+17}{Barroso, Clidaras, H\"{o}lzle. \textit{The Datacenter as a Computer.} 3rd ed., 2017.}
\mlsysref{Dean+13}{Dean \& Barroso. ``The Tail at Scale.'' CACM, 2013.}
\mlsysref{Crankshaw+17}{Crankshaw et al. ``Clipper: A Low-Latency Online Prediction Serving System.'' NSDI 2017.}
\mlsysref{MLPerf}{MLCommons. MLPerf Inference Benchmark Suite.}

\end{frame}

\begin{frame}{Next Lecture: ML Operations}
\note{[1 min] Forward hook. A single node is fragile. Models drift.
Deployments must roll out without downtime. Monitoring must detect silent
degradation. Ch14 scales from the single request to the full system lifecycle.}

\scriptsize
\begin{columns}[c]
  \begin{column}{0.30\textwidth}
    \centering
    \begin{mlsyscard}{computestroke}
    {\normalsize\bfseries CI/CD}\\[0.05cm]
    {\scriptsize Automated testing, canary deploys, rollback}
    \end{mlsyscard}
  \end{column}
  \begin{column}{0.30\textwidth}
    \centering
    \begin{mlsyscard}{datastroke}
    {\normalsize\bfseries Monitoring}\\[0.05cm]
    {\scriptsize Detect drift, accuracy decay, data quality}
    \end{mlsyscard}
  \end{column}
  \begin{column}{0.30\textwidth}
    \centering
    \begin{mlsyscard}{errorstroke}
    {\normalsize\bfseries Feature Stores}\\[0.05cm]
    {\scriptsize Consistent features across train and serve}
    \end{mlsyscard}
  \end{column}
\end{columns}

\vspace{0.2cm}
\centering
{\small From the single node to the \textbf{full system lifecycle}.}\\[0.05cm]
{\footnotesize How do we keep production ML systems running reliably through crashes, drift, and continuous updates?}

\end{frame}


% =============================================================================
% BACKUP SLIDES
% =============================================================================
\appendix

\begin{frame}{Backup: M/M/1 Queuing Worked Example}
\note{Use for additional depth or if students need alternative explanation.}
\small
Use if students struggle with the queuing theory calculations.\\small Service time: 5 ms ($\mu = 200$ req/s). Arrival rate: 140 req/s.\\ $\rho = 140/200 = 0.70$.\\ $W_q = \frac{0.70}{1 - 0.70} \times 5 = 11.7$ ms.\\ Total: $11.7 + 5 = 16.7$ ms.\\ p99 $\approx 4.6 \times 16.7 = 77$ ms.
\end{frame}

\begin{frame}{Backup: Continuous Batching Mechanics}
\note{Use for additional depth or if students need alternative explanation.}
\small
Use if students want deeper understanding of vLLM.\\small Static batching: all requests start and end together. Wasted compute on padding.\\ Continuous batching: when a request finishes, its slot is immediately filled.\\ PagedAttention: KV cache stored in fixed-size pages (like OS virtual memory).\\ Combined: 2--4$\times$ throughput improvement over naive batching.
\end{frame}

\begin{frame}{Backup: Cold Start --- The Problem}
\note{
% -- NARRATE: Expanded cold start timeline for deeper discussion
Walk through each phase: container pull, weight deserialization, graph
compilation (often the longest), memory pre-allocation, warmup inference.
The total can exceed 2 minutes for large LLMs.

% -- FLEX: [OPTIONAL] Pull if students want more detail on cold start
[OPTIONAL] Use when the main Cold Start slide felt too compressed.
}
\small
\textbf{Cold Start Timeline (Llama-3-70B on H100):}

\vspace{0.1cm}
\scriptsize
\renewcommand{\arraystretch}{1.1}
\begin{tabular}{@{}lrl@{}}
  \toprule
  \textbf{Phase} & \textbf{Time} & \textbf{Bottleneck} \\
  \midrule
  Container pull + startup & 10--30 s & Network + CPU \\
  Weight deserialization (140 GB) & 5--15 s & PCIe / NVLink \\
  Graph compilation (TensorRT) & 30--120 s & CPU \\
  KV cache pre-allocation & 2--5 s & GPU memory \\
  Warmup inference (compile kernels) & 3--10 s & GPU \\
  \midrule
  \textbf{Total} & \textbf{50--180 s} & \\
  \bottomrule
\end{tabular}

\vspace{0.1cm}
\mlsysalert{Key:}{Graph compilation is often the longest phase --- not weight loading. Serialized engines (TensorRT plan files) skip this step entirely.}
\end{frame}

\begin{frame}{Backup: Cold Start --- The Solutions}
\note{
% -- NARRATE: Mitigation strategies organized by effectiveness
Pre-warming keeps instances hot. Model caching avoids re-downloading.
Serialized engines skip compilation. Keep-alive prevents scale-to-zero.

% -- FLEX: [OPTIONAL] Pull alongside the Cold Start Problem slide
[OPTIONAL] Pair with the previous backup for a complete cold start module.
}
\small
\textbf{Cold Start Mitigations:}

\vspace{0.1cm}
\scriptsize
\renewcommand{\arraystretch}{1.1}
\begin{tabular}{@{}lp{4cm}r@{}}
  \toprule
  \textbf{Strategy} & \textbf{How It Works} & \textbf{Savings} \\
  \midrule
  \textbf{Pre-warming} & Keep $N$ instances always hot & 100\% (no cold start) \\
  \textbf{Model caching} & Cache weights on local NVMe SSD & 80\% of load time \\
  \textbf{Serialized engines} & Pre-compiled TensorRT plans & Skip 30--120 s compile \\
  \textbf{Keep-alive} & Never scale to zero & 100\% (no cold start) \\
  \textbf{Predictive scaling} & Scale up before traffic spike & Depends on prediction \\
  \bottomrule
\end{tabular}

\vspace{0.1cm}
\mlsysconcept{Trade-off:}{Pre-warming and keep-alive eliminate cold start but cost money during idle periods. Model caching is the best compromise for bursty traffic.}
\end{frame}

\begin{frame}{Backup: Tail at Scale --- Hedged Requests Worked Example}
\note{
% -- NARRATE: Concrete hedged request calculation
Service time: 5 ms (p50), 50 ms (p99). Without hedging: p99 = 50 ms.
With hedging (send to 2 replicas): p99 = $P(\text{both slow}) = 0.01^2
= 0.0001$. Effective p99 drops from 50 ms to approximately 10 ms.
Cost: 2x requests but 5x better tail.

% -- FLEX: [OPTIONAL] Pull for deeper Tail at Scale discussion
[OPTIONAL] Use if the war story generated strong student interest.
}
\small
Use if students want to calculate hedged request improvements.\\
\footnotesize
\textbf{Without hedging:}\\
$P(\text{slow}) = 0.01$. p99 = 50 ms.\\[0.1cm]
\textbf{With hedging (2 replicas):}\\
$P(\text{both slow}) = 0.01^2 = 0.0001$\\
Effective p99.99 = 50 ms. Effective p99 $\approx$ 10 ms.\\[0.1cm]
\textbf{Cost:} 2$\times$ requests. \textbf{Benefit:} 5$\times$ better tail latency.\\
Net: hedging is almost always worth it when tail latency matters.
\end{frame}

\end{document}
